% =====================================================================
%  Translation Without Meaning
%  Two-Column Relaxation to Semantic Quiescence in Finite
%  Weighted Graphs
%
%  A complete and rigorous formalisation. A translation is proper
%  when bidirectional cross-propagation between source and target
%  columns reaches a fixed point (quiescence); quiescence is identity
%  of meaning, certified without the translator ever knowing the
%  meaning. The floor replaces every imported uncertainty calculus.
%
%  Self-contained. Every object is a finite weighted graph or a
%  construction on one. Interpretation is confined to remarks.
% =====================================================================
\documentclass[11pt,a4paper]{article}

\usepackage[utf8]{inputenc}
\usepackage[T1]{fontenc}
\usepackage{lmodern}
\usepackage[a4paper,margin=2.4cm]{geometry}
\usepackage{amsmath,amssymb,amsthm,mathtools}
\usepackage{bm}
\usepackage{mathrsfs}
\usepackage{booktabs}
\usepackage{array}
\usepackage{enumitem}
\usepackage{microtype}
\usepackage{graphicx}
\usepackage{xcolor}
\usepackage[round,sort&compress]{natbib}
\usepackage[colorlinks=true,linkcolor=blue!45!black,
            citecolor=blue!45!black,urlcolor=blue!45!black]{hyperref}
\usepackage{cleveref}

% ---- Theorem environments ----
\newtheoremstyle{thmstyle}{\topsep}{\topsep}{\itshape}{}{\bfseries}{.}{.5em}{}
\theoremstyle{thmstyle}
\newtheorem{theorem}{Theorem}[section]
\newtheorem{lemma}[theorem]{Lemma}
\newtheorem{proposition}[theorem]{Proposition}
\newtheorem{corollary}[theorem]{Corollary}
\theoremstyle{definition}
\newtheorem{definition}[theorem]{Definition}
\newtheorem{axiom}{Premise}
\newtheorem{construction}[theorem]{Construction}
\newtheorem{example}[theorem]{Example}
\theoremstyle{remark}
\newtheorem{remark}[theorem]{Remark}
\newtheorem{notation}[theorem]{Notation}

\crefname{axiom}{Premise}{Premises}
\Crefname{axiom}{Premise}{Premises}

% ---- Shorthand ----
\newcommand{\R}{\mathbb{R}}
\newcommand{\N}{\mathbb{N}}
\newcommand{\Z}{\mathbb{Z}}
\newcommand{\Gph}{\Gamma}
\newcommand{\Lang}{L}
\newcommand{\V}{V}
\newcommand{\E}{E}
\newcommand{\wt}{w}
\newcommand{\med}{\mathsf{m}}          % medium
\newcommand{\flo}{\beta}               % floor
\newcommand{\cut}{\partial}            % cut
\newcommand{\co}{\complement}
\newcommand{\str}{\sigma}              % separation cost
\newcommand{\Elt}{\mathcal{X}}         % element set (catalogue)
\newcommand{\elt}{x}                   % an element
\newcommand{\hol}{\eta}                % holonomy
\newcommand{\Walk}{W}
\newcommand{\rec}{M}
\newcommand{\res}{\rho}                % residual / remaining uncertainty
\newcommand{\Rep}{\mathsf{R}}
\newcommand{\align}{a}                 % alignment
\newcommand{\abl}{\delta}              % ablation drop
\newcommand{\colA}{A}                  % source column
\newcommand{\colB}{B}                  % target column
\newcommand{\demand}{D}                % cross-demand
\newcommand{\relax}{\Phi}              % relaxation map
\newcommand{\abs}[1]{\lvert#1\rvert}
\newcommand{\eps}{\varepsilon}

\title{\textbf{Translation Without Meaning}\\[3pt]
\large Two-Column Relaxation to Semantic Quiescence\\
in Finite Weighted Graphs}

\author{%
Kundai Farai Sachikonye\\
\small Technical University of Munich\\
\small \texttt{kundai.sachikonye@wzw.tum.de}}

\date{\today}

\begin{document}
\maketitle

% =====================================================================
\begin{abstract}
\noindent
We formalise translation as a process that produces meaning-identical
output without the translator ever knowing the meaning of any word. The
construction rests on one inherited fact---that individuating a part of a
non-completable whole costs a strictly positive floor $\flo>0$---and uses
no imported probability, fuzzy, or Bayesian calculus: where prior treatments
reach for a borrowed uncertainty measure, the floor, which is derived and
parameter-free, does the work. We argue first that the object of translation
is not the \emph{content} of a text but the \emph{meaning} that propagates
from it as a cause, and that these differ: a precise translation reproduces
content with fidelity and thereby misses meaning, exactly as a precise mass
spectrum reproduces a reading without being the molecule's significance.
Since meaning is an invariant that cannot be extracted---accessing it would
require the accessor to \emph{be} the item---we do not try to know it; we
detect when two sides \emph{share} it, by the absence of any propagable
difference between them. We model each language as a contact graph and each
side of a translation as a column; the translator runs a \emph{bidirectional
relaxation}: perturb a column, reverse-translate the perturbation into the
other, propagate any demanded change back, and iterate. We prove:
\textbf{(Q)} the relaxation is monotone and its demands are floor-bounded, so
it is a contraction toward a fixed point or it does not halt; \textbf{(I)}
\emph{at quiescence the two sides have the same meaning}---necessarily,
because any difference of meaning is itself a propagable demand and
quiescence is the absence of all demands, so quiescence does not approximate
semantic identity but constitutes it, certified without the meaning ever
being known; \textbf{(F)} quiescence is reached at positive residual
$\res\ge\flo$ (same meaning, irreducibly different form), never at zero
(exactness, identity of form, is forbidden); \textbf{(U)} non-halting is the
proof of deep untranslatability: no fixed point exists, the two carvings
admit no mutual semantic equilibrium, and the honest output is to decline;
and \textbf{(S)} a separating experiment---giving the three kinds of
translator a grammatically wrong statement---shows the human depends on
meaning (translates it by routing around the broken form), the
content-substituter depends on form (mangles it), and our method depends on
neither but on the existence of a fixed point (translates it iff the
relaxation still reaches quiescence, and declines otherwise where the human
infers and the substituter hallucinates). The elements that the relaxation
holds fixed are found by \emph{ablation}, of emergent grain; faithfulness is
zero back-holonomy; the perturbation suite measures the framework's own
floor, opacity, mobility, and holonomy. This is translation without meaning:
meaning-identical output by a meaning-blind, comprehension-free mechanism,
sufficient exactly where a fixed point exists. Every claim is a construction
on finite weighted graphs; interpretation is confined to remarks.
\end{abstract}

\noindent\textbf{Keywords:} translation; finite weighted graph; meaning as
region; content versus meaning; bidirectional relaxation; semantic
quiescence; fixed point; contact floor; ablation; cycle holonomy;
faithfulness; deep untranslatability; comprehension-free translation.

\tableofcontents

% =====================================================================
\part{The Object of Translation}
% =====================================================================

% =====================================================================
\section{Introduction: translation without meaning}
\label{sec:intro}
% =====================================================================

\subsection{The thesis}

We formalise a translator that never knows what any word means and
nonetheless produces output identical in meaning to the source. The two
claims are not in tension once one sees what meaning is and how identity of
meaning can be detected without knowing it.

Three facts organise the paper.

\emph{First}, a text exists only because uncertainty propagated. There is no
statement to translate unless something not-yet-known was carried forward
into marks; a text is not data but a \emph{cause}---a propagation of
residue---and the meaning of a text is what that cause emits, not the marks
themselves. A conversation runs by propagation: its participants are already
uncertain, hold no ledger of points to satisfy in order, and do not know
in advance what they will get out of it. What they extract is irreducible---
not recoverable from the inputs, or there would have been no need to
converse.

\emph{Second}, meaning is not content, and not intent. ``How are you?''
in British usage does not mean that the speaker wishes to be told the state
of one's life; its content is a health-inquiry, its meaning is a greeting,
and the two have come apart. Asking a person's age has the content
``what is your age'' but may carry the meaning ``can this person vote, or
drive''---which may not even be what one wanted. Meaning is downstream of the
cause and is generally \emph{not} the literal content nor the speaker's
intent.

\emph{Third}, meaning cannot be extracted. A cell is individuated only by
what it is not, against a non-completable whole; to extract its meaning---to
complete the comparison that fixes it---one would have to occupy its entire
negation-structure, that is, to \emph{be} it, which is no longer a
comparison. So no translator, human or machine, reads meaning off a word as
a value. What a translator can do is something weaker and sufficient: detect
when two sides \emph{share} a meaning, by the absence of any difference that
would propagate between them.

From these: a translator that does not know meaning can still translate, by
relaxing the two sides until no difference can propagate; and when the
relaxation halts, the two sides mean the same thing---not approximately, but
exactly, because the only thing that could distinguish their meanings is a
difference, and a difference is precisely what would keep the relaxation
running.

\subsection{Why precision is the wrong target}

\begin{remark}[A precise translation is a precise reading of content]
\label{rem:massspec}
Optimising a translation for precision optimises fidelity to the source
\emph{content}---the literal marks and their cognates. By the second fact
above, content is exactly what the meaning is not. A precise translation
therefore reproduces the non-meaning with high fidelity, as a precise mass
spectrum reproduces a reading that is not the molecule's significance: the
spectrum is an accurate \emph{contact}, and what one wants---what the
molecule is, or does, or lets one deduce---is the propagation, not the
contact. ``How are you?'' rendered precisely yields its health-inquiry
cognate and misses the greeting. Precision is fidelity to the wrong object;
the right object is the meaning, which is the propagation of the source as a
cause. This remark motivates the criterion of \Cref{sec:relaxation}, on which
the theorems depend; no theorem depends on the mass-spectrum reading itself.
\end{remark}

\subsection{The single inherited fact}

\begin{axiom}[Contact floor]
\label{ax:floor}
A \emph{contact graph} is a finite, connected, weighted graph
$\Gph=(\V,\E,\wt)$ with a distinguished \emph{medium} vertex $\med$ adjacent
to all others, and with a strictly positive floor on edge weight: there is
$\flo>0$ with $\wt(e)\ge\flo$ for all $e\in\E$. Equivalently every cut
separating a proper part from the rest has weight $\ge\flo$. The floor is
derived, not posited: to complete the comparison that individuates a part,
the comparer would have to bring the entire complement to bear, which---since
the part is fixed only as the complement of everything it is not---would
require the comparer to occupy the part's own negation-structure, i.e.\ to
\emph{be} the part, which is no longer a comparison; so the comparison never
completes and an irreducible residual $\ge\flo$ always remains. Here we use
only $\flo>0$.
\end{axiom}

\begin{remark}[The floor replaces the borrowed calculus]
\label{rem:floor-replaces}
The role usually played by a probability, a fuzzy membership, or a Bayesian
posterior---a number in $[0,1]$ recording how sure or how much---is played
throughout by \emph{distance to the floor}: a meaning is resolved to within
the floor, never to a point; an irreducible residual $\res\ge\flo$ always
remains; a translation relaxes toward the floor. The floor is not fitted---it
is the trace the non-completable whole leaves on each finite part---so no
membership function is tuned and no prior is set. Where an earlier
specification recorded a field ``remaining\_uncertainty: a probability,'' we
write ``the residual $\res\ge\flo$,'' derived. No theorem depends on the
comparison.
\end{remark}

% =====================================================================
\section{Languages, meanings, and the impossibility of cognate-substitution}
\label{sec:language}
% =====================================================================

\begin{definition}[Meaning-graph of a language]
\label{def:meaning-graph}
A \emph{language} $\Lang$ is modelled by a contact graph
$\Gph_\Lang=(\V_\Lang,\E_\Lang,\wt_\Lang;\med_\Lang)$ whose vertices are
\emph{meaning-cells} the language holds distinct, whose edge $\{u,v\}$ is a
\emph{contact} of cost $\wt_\Lang(\{u,v\})\ge\flo$ (the cost of telling $u$
from $v$), and whose medium $\med_\Lang$ is ``the rest of what the language
can mean,'' against which each cell is individuated.
\end{definition}

\begin{definition}[Meaning is a region; separation cost]
\label{def:region}
The \emph{separation cost} of a cell $v$ is $\str_\Lang(v)=
\min_{S\ni v,\ \med\notin S}\wt(\cut(S))\ge\flo$, with minimising
\emph{resting cut} $\cut(S^{\ast}(v))$. The \emph{meaning} of $v$ is this
resting cut---the region of contrasts by which $v$ is told from the rest,
of weight $\ge\flo$. It is never a point: a point would be a zero-weight cut,
forbidden by \Cref{ax:floor}.
\end{definition}

\begin{definition}[Content versus meaning]
\label{def:content}
The \emph{content} of a source unit is its literal form---its marks, tokens,
and their dictionary cognates. The \emph{meaning} is its resting cut
(\Cref{def:region}). The content is a label on the cell; the meaning is the
cell's individuation against the whole. They coincide only when the literal
form is itself the operative individuation, which is not the general case
(\Cref{ex:howareyou}).
\end{definition}

\begin{example}[Content and meaning come apart]
\label{ex:howareyou}
``How are you?'' has content \emph{health-inquiry} and, in British greeting
usage, meaning \emph{greeting}: its resting cut individuates it against
other greetings, not against other health-inquiries. A translator preserving
content yields the health-inquiry cognate; a translator preserving meaning
yields the target's greeting. The two land on different target cells.
\end{example}

\begin{theorem}[Cognate-substitution cannot reach meaning]
\label{thm:no-cognate}
Let two languages $\Lang,\Lang'$ have meaning-graphs with distinct weights.
Then a content map---one sending each source token to its target cognate---
does not in general preserve meaning: there exist source units whose cognate
image has a different resting cut from the unit's meaning. Two structural
obstructions guarantee this:
\begin{enumerate}[label=\textup{(\roman*)},leftmargin=2.2em]
\item \textbf{Different carvings.} $\Lang$ and $\Lang'$ individuate the space
of senses differently, so a source cell's meaning need not be any single
target cell (\Cref{ex:colour}).
\item \textbf{Content $\neq$ meaning.} Even within one carving, a unit's
meaning is its resting cut, not its content; the cognate preserves content,
which is not the meaning (\Cref{ex:howareyou}).
\end{enumerate}
\end{theorem}

\begin{proof}
\textup{(i)} Let $\Lang$ resolve one cell $\elt$ where $\Lang'$ resolves two,
$\elt'_1,\elt'_2$ (\Cref{ex:colour}); the content map sends $\elt$ to one
cognate, whose resting cut individuates it against the \emph{other} target
cell, differing from $\elt$'s resting cut, which has no such internal
contrast. \textup{(ii)} In \Cref{ex:howareyou} the cognate preserves the
health-inquiry content; its resting cut individuates against health-inquiries,
while the unit's meaning individuates against greetings; the two cuts differ.
In both cases the cognate image's meaning $\neq$ the source meaning.
\end{proof}

\begin{example}[The colour terms]
\label{ex:colour}
English resolves one cell, \emph{blue}, where Russian resolves two,
\emph{goluboj} and \emph{sinij}. These are different individuations of one
space, not a finer reading of shared points; no single Russian cell covers
\emph{blue}, only the composite. Substitution has nothing to substitute.
\end{example}

% =====================================================================
\section{Elements by ablation, with emergent grain}
\label{sec:elements}
% =====================================================================

We need the items a unit cannot do without---its \emph{elements}---because
the relaxation will hold these fixed and move only the rest. Elements are
defined intrinsically and operationally, by ablation, computable on a single
unit, with no whole-language catalogue built in advance.

\begin{definition}[Ablation drop; replaceability; element; residue]
\label{def:ablation}
Let $U$ be a meaning-unit (a connected sub-graph) with separation cost
$\str_\Lang(U)$, and $i$ an item of $U$. The \emph{ablation drop} of $i$ is
\[
\abl(i\mid U)\ :=\ \str_\Lang(U)\ -\ \str_\Lang(U\setminus i).
\]
$i$ is \emph{replaceable} to the degree $\abl(i\mid U)$ is small (its removal
barely lowers the unit's individuation; some other item would serve) and
\emph{unreplaceable} to the degree $\abl(i\mid U)$ is large (its removal
collapses the unit; none serves). The \emph{elements} of $U$ are its
unreplaceable items (maximal $\abl$); the \emph{residue} is the replaceable
remainder.
\end{definition}

\begin{theorem}[Unreplaceable $=$ irreducible $=$ at the floor]
\label{thm:element-floor}
$i$ is an element of $U$ iff it is irreducible in $U$ (not splittable into,
or substitutable by, a cheaper sub-individuation without raising $U$'s total
cost), iff its contribution to $U$'s resting cut sits at the floor---an
irreducible $\ge\flo$ of the unit's distinctness no replacement supplies.
\end{theorem}

\begin{proof}
If $i$ were reducible, a cheaper substitute would recover its contribution,
so removing $i$ would not collapse $U$ and $\abl(i\mid U)$ would be small:
$i$ replaceable. Conversely if no substitute reproduces $i$'s contribution,
removing $i$ lowers $U$'s individuation by the floor amount $i$ uniquely
carries, so $\abl(i\mid U)\ge\flo$: $i$ unreplaceable. The amount $i$ carries
is $\ge\flo$ (\Cref{ax:floor}) and equals $\abl(i\mid U)$ at the
reduction-terminal stage.
\end{proof}

\begin{theorem}[The catalogue exists; the grain is an output]
\label{thm:catalogue}
Every unit $U$ has a finite, nonempty element-set $\Elt_\Lang(U)$. The
\emph{grain} of its elements (sub-word morpheme, single word, multi-word
group) is determined by ablation and not stipulated: an element is whatever
item, at whatever span, whose ablation collapses the unit.
\end{theorem}

\begin{proof}
Ablation ranks finitely many items of a finite unit; the maximal-$\abl$ set
exists, is finite, and is nonempty because $U$'s positive individuation
$\ge\flo$ must be carried by some item. $\abl(i\mid U)$ is defined for $i$ a
morpheme, a word, or a group equally, so the element is whichever span
maximises the drop; the span is read off, not imposed.
\end{proof}

\begin{remark}[Idioms, morphemes, and ``binds to'']
\label{rem:grain}
\emph{Kicked the bucket} is one multi-word element: ablating any single word
destroys the idiomatic individuation (large group-$\abl$) while the literal
sub-senses do not compose to it, so the group is irreducible. \emph{Binds to}
is one element when the pair carries the relation though each word's
individual $\abl$ is small. A bound tense or case morpheme of irreducible
grammatical sense is a sub-word element. This is the \emph{fuzzy unit} of an
earlier specification: a definitely-belonging core (the element, high $\abl$)
with graded-membership boundary material (the residue, low $\abl$), scale-
fluid. No theorem depends on the reading.
\end{remark}

\subsection{Position fixes the slot}

\begin{definition}[Slot; substitution-set]
\label{def:slot}
In $U$, the \emph{slot} of $i$ is its positional-semantic role (agent,
action, recipient, manner, time, \dots), fixed by $i$'s position and the
grammar of $\Lang$. The \emph{substitution-set} $\mathrm{Sub}(i\mid U)$ is
the set of items the grammar admits in that slot.
\end{definition}

\begin{theorem}[Replaceability is the cut against the slot]
\label{thm:slot}
$i$'s replaceability equals the cheapness of individuating $i$ from
$\mathrm{Sub}(i\mid U)$ in the context of $U$: $i$ is replaceable iff
$\str_\Lang(i,\mathrm{Sub}(i\mid U)\mid U)$ is small, and an element iff that
cut is maximal. This agrees with ablation: an item barely told from its
slot-alternatives ablates cheaply; an item far from all alternatives ablates
catastrophically.
\end{theorem}

\begin{proof}
A small slot-cut means some $j\in\mathrm{Sub}(i\mid U)$ lies within the floor
of $i$, so $j$ recovers $i$'s contribution: $i$ replaceable, small $\abl$. A
maximal slot-cut means every alternative lies a fixed margin above the floor,
so none recovers $U$'s individuation: $i$ an element, large $\abl$. The two
orderings coincide.
\end{proof}

\begin{remark}[Position is meaning; the reversal]
\label{rem:position}
\Cref{thm:slot} is ``the location of a word is the point behind its
meaning,'' made precise. ``The bank approved the loan'' versus ``the loan
approved the bank'': the same tokens, different slots, hence different
substitution-sets and different individuations. A bag-of-tokens map is blind
to exactly the slot structure that carries the sense. No theorem depends on
the reading.
\end{remark}

% =====================================================================
\part{The Relaxation and Its Fixed Point}
% =====================================================================

% =====================================================================
\section{Termination of units; the naive seed}
\label{sec:setup}
% =====================================================================

\begin{definition}[Terminated unit]
\label{def:terminated}
A source unit is \emph{terminated} if it is a bounded cell with a definite
resting cut, hence a separation cost $\ge\flo$. It is \emph{open} if it has
no resting cut (the minimising cut is not attained).
\end{definition}

\begin{theorem}[Only terminated units enter the relaxation]
\label{thm:terminate}
A unit can be cross-propagated (\Cref{def:relaxation}) only if terminated.
An open unit has no defined cut against the medium, hence no defined
alignment to any target cell, hence no demand to propagate.
\end{theorem}

\begin{proof}
The cross-demand of \Cref{def:demand} is defined through alignment, a
minimum-cut weight requiring a defined cut against the medium---termination.
An open unit's alignment is undefined; no demand is computable.
\end{proof}

\begin{construction}[Perfective termination]
\label{con:perfective}
Cast each source clause in a completed (perfective) aspect---bounded
(``has done'') rather than open (``was doing'')---giving every unit a
completion boundary, the resting cut the relaxation requires; restore natural
aspect on output. Some termination must be applied; the perfect is one the
grammar supplies. \textit{(No theorem depends on the choice.)}
\end{construction}

\begin{construction}[The two columns; the naive seed]
\label{con:columns}
Place the terminated source as column $\colA$. Run a naive cognate
translation---each item to its dictionary cognate---as column $\colB$,
the \emph{seed}: aligned to $\colA$, usually partly wrong, but structured.
The seed is the initial (off-target) state the relaxation will correct. In
each column, ablation-rank the items (\Cref{def:ablation}); the high-$\abl$
items are the column's \emph{anchors} (elements, at the floor), held fixed;
the residue is what the relaxation moves.
\end{construction}

% =====================================================================
\section{The bidirectional relaxation}
\label{sec:relaxation}
% =====================================================================

\subsection{Cross-demand}

\begin{definition}[Meaning-medium and alignment]
\label{def:alignment}
Both languages individuate one space of senses; embed their graphs in the
common \emph{meaning-medium} $\Gph$---the shared, uncarved whole, not a third
language---joined by contacts wherever a cell of one is individuable against a
cell of the other. For cells $u,v$ of either column the \emph{alignment} is
$\align(u,v)=\min_{S\ni u,\ v\notin S}\wt(\cut(S))\ge\flo$, normalised into
$(0,1]$.
\end{definition}

\begin{remark}[The medium is the shared world, not a posit]
\label{rem:medium}
The medium is not an extra assumption: both languages are finite
individuations of the \emph{same} non-completable whole, so they share it as
two finite contact graphs of one world share a medium. They differ
(\Cref{thm:no-cognate}) in their carvings, not in the space carved. What is
empirical is which cross-contacts hold; that is the work the method
organises. No theorem depends on the reading.
\end{remark}

\begin{definition}[Reverse-translation; cross-demand]
\label{def:demand}
Let $\relax_{\colA\to\colB}$ be the operation that takes the current state of
$\colA$, reverse-translates it into $\colB$'s language by residue propagation
(\Cref{def:propagation}), and returns, for each residue item of $\colB$, the
change required to make $\colB$ consistent with $\colA$'s reverse-image. The
\emph{cross-demand} $\demand_{\colA\to\colB}$ is the total such required
change, measured as the summed alignment gap between $\colB$'s current
residue cells and the cells $\colA$'s reverse-image places there:
\[
\demand_{\colA\to\colB}\ :=\ \sum_{i\in\mathrm{residue}(\colB)}
\bigl(\align(\colB_i,\ \relax_{\colA\to\colB}(\colA)_i)-\flo\bigr)\ \ge\ 0.
\]
$\demand_{\colB\to\colA}$ is defined symmetrically. A column is
\emph{satisfied from the other} when that direction's cross-demand is zero
(every residue cell already at floor distance from what the other side
demands).
\end{definition}

\begin{definition}[Residue propagation]
\label{def:propagation}
A \emph{residue propagation} reverse-translating column $X$ into the other
language is a walk in the meaning-medium from each of $X$'s residue cells to
a cell of the other language, holding $X$'s anchors fixed; each step commits
residue $\ge\flo$ and the cause of the next step is the residual deposited by
the last; intermediate cells may belong to either language or to neither.
\end{definition}

\subsection{The relaxation step and its fixed point}

\begin{definition}[Relaxation; quiescence; proper translation]
\label{def:relaxation}
The \emph{relaxation} alternates: apply $\relax_{\colA\to\colB}$ (update
$\colB$'s residue to satisfy $\colA$, then re-ablate), then
$\relax_{\colB\to\colA}$ (update $\colA$'s residue to satisfy $\colB$), and
iterate. The pair $(\colA,\colB)$ is \emph{quiescent}---a \emph{fixed point}
of the relaxation---when both cross-demands vanish:
$\demand_{\colA\to\colB}=\demand_{\colB\to\colA}=0$, so that a change
propagated from either side requires no change on the other. A translation is
\emph{proper} iff it is quiescent.
\end{definition}

\begin{theorem}[The relaxation is monotone and floor-bounded]
\label{thm:monotone}
Along the relaxation the total committed residue is non-decreasing (each
update commits cuts $\ge\flo$ and none is withdrawn---withdrawal is itself a
new committed cut), and the residue is bounded below at every step by $\flo$
per maintained contact. Hence the sequence of committed-residue totals is
monotone non-decreasing, and the sequence of cross-demands is non-increasing
once anchors have stabilised.
\end{theorem}

\begin{proof}
Each update is a residue propagation (\Cref{def:propagation}) committing cuts
of weight $\ge\flo$; by the monotone, non-returning character of committed
cuts (\Cref{ax:floor} forbids zero-cost cuts, and un-committing redraws a
boundary, itself a committed cut), the committed total never decreases. Once
the anchors (the high-$\abl$ elements) no longer change---which occurs after
finitely many updates, since each column has finitely many items and an
element re-classified by an update can be re-classified at most until its
$\abl$ ordering is stable---each subsequent update only moves residue toward
satisfying the other side's demand, which by definition reduces that
direction's cross-demand or leaves it unchanged. Hence cross-demands are
non-increasing thereafter.
\end{proof}

\begin{theorem}[Dichotomy: quiescence or non-halting]
\label{thm:dichotomy}
The relaxation either reaches quiescence in finitely many steps, or it does
not halt. If the cross-demands reach zero, $(\colA,\colB)$ is a fixed point
(proper). If they remain bounded below by a positive constant across
unboundedly many post-stabilisation updates, no fixed point is reached and
the relaxation does not halt.
\end{theorem}

\begin{proof}
After anchor stabilisation (\Cref{thm:monotone}) the cross-demands are
non-increasing and bounded below by $0$. A non-increasing sequence bounded
below converges. If its limit is $0$ and the state space is finite (finitely
many residue assignments over finitely many cells), the infimum is attained
at a finite step: quiescence. If its limit is a positive constant, the
demands never vanish: each update still requires a change on the other side,
and the relaxation does not halt. These are the only two cases.
\end{proof}

% =====================================================================
\section{Quiescence is identity of meaning}
\label{sec:identity}
% =====================================================================

This is the central result: when the relaxation halts, the two sides mean the
same thing---not approximately, but exactly---and this is certified without
the meaning ever being known.

\begin{theorem}[Quiescence constitutes semantic identity]
\label{thm:quiescence-identity}
If $(\colA,\colB)$ is quiescent (\Cref{def:relaxation}), then $\colA$ and
$\colB$ have the same meaning: there is no difference of meaning between them.
Conversely, if $\colA$ and $\colB$ have a difference of meaning, the relaxation
is not quiescent. Hence
\[
\text{quiescence}\quad\Longleftrightarrow\quad\text{identity of meaning},
\]
and the equivalence holds without any party knowing what the meaning is.
\end{theorem}

\begin{proof}
\textup{(Quiescent $\Rightarrow$ identical.)} Suppose, for contradiction,
$(\colA,\colB)$ is quiescent yet the two differ in meaning. A difference of
meaning is a difference in resting cut: there is a contrast individuating
$\colA$ that does not individuate $\colB$, or conversely. Reverse-translate
$\colA$ into $\colB$'s language (\Cref{def:demand}); where $\colA$'s meaning
carries a contrast $\colB$ lacks, the reverse-image places a cell that
$\colB$'s current residue does not match, so $\align(\colB_i,
\relax_{\colA\to\colB}(\colA)_i)>\flo$ for that item, whence
$\demand_{\colA\to\colB}>0$. This contradicts quiescence
($\demand_{\colA\to\colB}=0$). Hence no difference of meaning can exist at
quiescence. \textup{(Difference $\Rightarrow$ not quiescent.)} The same
construction shows any meaning-difference yields a positive cross-demand,
which is non-quiescence. \textup{(Without knowing the meaning.)} The argument
refers only to cross-demands, which are alignment gaps---minimum-cut weights
between cells---computed without resolving any cell's meaning to a value; the
meaning (the resting cut, the unextractable invariant) is never read. Identity
is certified by the \emph{absence of a propagable difference}, not by
inspection of the meaning.
\end{proof}

\begin{corollary}[The translator detects identity it cannot read]
\label{cor:blind}
The translator never knows the meaning of any cell---meaning is the
unextractable invariant (\Cref{ax:floor})---yet, by relaxing to quiescence,
it produces a target whose meaning is identical to the source's. It achieves
semantic identity it cannot inspect, because identity is no-difference and
no-difference is testable blind.
\end{corollary}

\begin{proof}
Immediate from \Cref{thm:quiescence-identity}: identity is equivalent to the
vanishing of cross-demands, and cross-demands are computed without reading any
meaning.
\end{proof}

\begin{theorem}[Form survives: quiescence at positive residual]
\label{thm:form-residual}
A proper translation is quiescent at residual $\res\ge\flo>0$, never at
$\res=0$. The two sides have identical meaning and irreducibly different
form: $\res$ is the residue that does not cross-propagate at
quiescence---the difference of expression the two languages cannot erase
while remaining two languages.
\end{theorem}

\begin{proof}
By \Cref{ax:floor} every maintained contact carries weight $\ge\flo$; the two
columns are distinct individuations (\Cref{thm:no-cognate}), so even at
identity of meaning each retains its own form-cells, individuated at cost
$\ge\flo$ that the other cannot absorb without becoming the same language.
Thus the residual at quiescence is $\ge\flo>0$. It is not a residue of
meaning (cross-demands are zero) but of form. $\res=0$ would be a zero-cost
identity of expression, forbidden.
\end{proof}

\begin{remark}[Why this is the right criterion, not precision]
\label{rem:right-criterion}
A precision-seeking translator halts when the target matches the source
\emph{content}; the relaxation halts when no \emph{meaning}-difference can
propagate. These differ. ``How are you?'' relaxes thus: reverse-translate the
target greeting back into the source language; it does not demand the
health-inquiry content (a greeting is not a health-inquiry), so propagating
back demands no change; the pair is quiescent at the greeting, not at the
health-inquiry cognate. The precision-seeker halts early, at content; the
relaxation halts at meaning. Precision is fidelity to form
(\Cref{rem:massspec}); quiescence is identity of meaning
(\Cref{thm:quiescence-identity}). No theorem depends on the example.
\end{remark}

% =====================================================================
\section{Non-halting is deep untranslatability}
\label{sec:untranslatable}
% =====================================================================

\begin{theorem}[Two kinds of untranslatability]
\label{thm:two-kinds}
There are exactly two ways a source unit fails to translate, and the
relaxation distinguishes them:
\begin{enumerate}[label=\textup{(\roman*)},leftmargin=2.2em]
\item \textbf{Shallow (catalogue) untranslatability.} The relaxation
\emph{halts}, but no single target element is at floor distance; the unit is
rendered by a composite (gloss, circumlocution) or, if no composite reaches
floor distance, by a borrowing. This is a difference of \emph{carving},
resolved at quiescence by a multi-cell target.
\item \textbf{Deep untranslatability.} The relaxation \emph{does not halt}:
the cross-demands never vanish (\Cref{thm:dichotomy}), so no fixed point
exists and the two carvings admit no mutual semantic equilibrium. The
meanings \emph{cannot be made identical} across the two languages; the
perpetual demand is the proof.
\end{enumerate}
\end{theorem}

\begin{proof}
\textup{(i)} If the relaxation halts (quiescent), the meanings are identical
(\Cref{thm:quiescence-identity}); identity may be realised by a single target
element at floor distance, or, when the target carves more finely or coarsely
(\Cref{thm:no-cognate}(i)), by the minimal composite of target elements whose
joint cell reaches floor distance, or by a borrowing when no composite does.
The trichotomy direct/compose/borrow is decided by computing single- and
composite-element alignments at the fixed point. \textup{(ii)} If the
relaxation does not halt, by \Cref{thm:dichotomy} the cross-demands are
bounded below by a positive constant; by \Cref{thm:quiescence-identity} a
nonzero cross-demand is a difference of meaning; so a difference of meaning
persists under every update---no assignment of target cells makes the meanings
identical. The meanings cannot be equalised: deep untranslatability.
\end{proof}

\begin{corollary}[The honest output: decline, not hallucinate]
\label{cor:decline}
On a deeply untranslatable unit the relaxation does not reach quiescence, and
the correct output is to \emph{decline}---to report that no proper
translation exists---rather than to emit a non-quiescent (meaning-divergent)
rendering. The method does not fabricate where no fixed point exists.
\end{corollary}

\begin{proof}
By \Cref{thm:two-kinds}(ii) non-halting means no quiescent pair, hence no
proper translation (\Cref{def:relaxation}); any emitted rendering would be
non-quiescent, i.e.\ meaning-divergent (\Cref{thm:quiescence-identity}).
Declining is the only output consistent with properness.
\end{proof}

\begin{remark}[Borrowing grows the catalogue irreversibly]
\label{rem:borrow}
A borrowing commits a new cut in the target graph, growing its catalogue by
one; by non-return this is irreversible---a loanword, once in, is a new
element. Loanwords accrete and do not un-accrete because each is a committed
individuation. No theorem depends on the reading.
\end{remark}

% =====================================================================
\part{Evidence, Faithfulness, and the Pipeline}
% =====================================================================

% =====================================================================
\section{The separating experiment: three translators, one test}
\label{sec:separating}
% =====================================================================

We give an experiment that dissociates the two things a translator might
depend on---meaning and form---and thereby identifies which each of three
methods depends on. A dissociating input is a \emph{grammatically wrong}
statement: one whose form is broken but whose meaning may still be
recoverable.

\begin{theorem}[The grammatically-wrong dissociation]
\label{thm:separating}
Present a grammatically wrong statement \textup{(}broken form, possibly
recoverable meaning\textup{)} to three translators. Then:
\begin{enumerate}[label=\textup{(\roman*)},leftmargin=2.2em]
\item A \textbf{meaning-dependent} translator (the human, who derives meaning
from context) translates it: it routes around the broken form to the
meaning. Success on broken form proves dependence on meaning.
\item A \textbf{form-dependent} translator (content-substitution: cognate
lookup, statistical smoothing) mangles or fails it: it has no well-formed
structure to substitute over. Failure on broken form proves dependence on
form.
\item Our \textbf{fixed-point-dependent} method translates it \emph{iff the
relaxation still reaches quiescence}---iff the broken statement still admits a
recoverable, propagable meaning---and \emph{declines otherwise}
(\Cref{cor:decline}). It depends on neither meaning (which it never knows) nor
clean form (which it never substitutes), but on the existence of a fixed
point.
\end{enumerate}
\end{theorem}

\begin{proof}
\textup{(i)} A method that succeeds when the form is broken cannot be using
the form as its input; the only remaining carrier of the rendering is the
meaning; hence the method depends on meaning. \textup{(ii)} A method that
fails when the form is broken is using the form; hence it depends on form.
\textup{(iii)} Our method's machinery (slot $\to$ substitution-set $\to$
ablation $\to$ cross-demand) needs slots, which broken grammar damages; where
enough slot-structure survives that the relaxation can still run to
quiescence, the meaning is propagably present and the method translates it
(blind, \Cref{cor:blind}); where the damage destroys the propagable structure,
no fixed point exists (\Cref{thm:two-kinds}(ii)) and the method declines
(\Cref{cor:decline}). It depends on quiescence-reachability---the existence of
a fixed point---which is the recoverable presence of meaning, not its
comprehension and not the form.
\end{proof}

\begin{corollary}[Located between human and substituter; honest at the boundary]
\label{cor:located}
Our method reaches the meaning-dependent translator's \emph{result} (identical
meaning) without its \emph{comprehension}, by relaxation; and it never emits
the form-dependent translator's mangling, because it does not substitute. At
the boundary---inputs whose meaning is recoverable only by inference beyond
any propagable structure---the human \emph{infers}, the substituter
\emph{hallucinates}, and our method \emph{declines}. The method is sufficient
exactly on inputs admitting a fixed point.
\end{corollary}

\begin{proof}
Reaching the result without comprehension is \Cref{cor:blind}; not mangling is
that the method emits only quiescent (meaning-identical) output or declines
(\Cref{cor:decline}); the boundary behaviour is the contrast of
\Cref{thm:separating}(i)--(iii) at the edge of fixed-point existence.
\end{proof}

% =====================================================================
\section{Faithfulness and the perturbation suite}
\label{sec:faithful}
% =====================================================================

\subsection{Faithfulness is zero back-holonomy}

\begin{definition}[Meaning-transport; back-cycle; holonomy]
\label{def:holonomy}
Let $g$ assign to each contact an element of an abelian group---the
sense-shift the contact carries---with $g(uv)=-g(vu)$. For a forward
propagation $\Walk$ of $\colA$ to $\colB$ and a back-propagation $\Walk^{-}$,
the \emph{back-cycle} holonomy is $\hol=\sum_{e\in\Walk}g(e)+
\sum_{e\in\Walk^{-}}g(e)$.
\end{definition}

\begin{theorem}[Faithfulness is back-cycle closure; quiescence implies it]
\label{thm:faithful}
A translation is \emph{faithful} iff some back-propagation has $\hol=0$
(translating across and back returns to the source, to within the floor). At
quiescence the back-cycle closes: $\demand_{\colB\to\colA}=0$ means
reverse-translating $\colB$ demands no change in $\colA$, which is
$\hol=0$. Hence a proper translation is faithful, and an unfaithful one is
non-quiescent.
\end{theorem}

\begin{proof}
A group-valued transport is path-independent between two endpoints iff every
cycle between them has zero holonomy; the round trip $\colA\to\colB\to\colA$
returns to the source iff its cycle closes. Quiescence is
$\demand_{\colB\to\colA}=0$ (\Cref{def:relaxation}), i.e.\ the back direction
demands no change---exactly that the round trip returns to $\colA$, i.e.\
$\hol=0$. So proper $\Rightarrow$ faithful; $\hol\neq0$ for all back-routes
$\Rightarrow$ a standing demand $\Rightarrow$ non-quiescent.
\end{proof}

\begin{theorem}[The false friend is a bridge]
\label{thm:falsefriend}
A \emph{false friend}---a target cell whose surface matches the source but
whose meaning does not---has \emph{endpoints} that appear aligned (each a
real, locally plausible cell) but a \emph{back-cycle} with nonzero holonomy
through an intermediate sense. It passes an endpoint check and fails a route
check; it is non-quiescent (\Cref{thm:faithful}) and detectable only by
auditing the route, not the endpoints.
\end{theorem}

\begin{proof}
The endpoints are real, locally matched cells, so an endpoint invariant
registers a plausible pairing. The meaning is carried by the interior: the
back-cycle through the intermediate sense has $\hol\neq0$ because the
candidate's actual cell is not the source's; by \Cref{thm:faithful} the pair
is non-quiescent, with the standing demand on the interior. No endpoint-only
check detects it; back-translating through the meaning does. This is the
bridge: closing endpoints, non-closing route.
\end{proof}

\subsection{The perturbation suite measures the floor}

\begin{definition}[Perturbation suite; stability]
\label{def:perturb}
For a translation of $U$, the \emph{perturbation suite} applies four families
and records the change in landing: \textup{(1)} ablation (remove each item);
\textup{(2)} in-grammar rearrangement; \textup{(3)} synonym substitution;
\textup{(4)} negation. The \emph{stability} is $1-\bar{\Delta}/\Delta_0$,
mean landing-change over base landing.
\end{definition}

\begin{theorem}[Each perturbation measures a proved structure]
\label{thm:perturb}
The four families measure four structures of the formalisation:
\textup{(i)} ablation measures the element/residue split (high-$\abl$ items
are the elements, \Cref{thm:element-floor}---the ablation profile \emph{is}
the element-set); \textup{(ii)} in-grammar rearrangement measures path opacity
(invariance under legal reorder is endpoint-determination); \textup{(iii)}
synonym substitution measures representation mobility (invariance under a
same-slot synonym marks a landing on an element, not a representation
artifact); \textup{(iv)} negation measures holonomy (a faithful landing must
invert correctly under negation, i.e.\ the negation-cycle closes,
\Cref{thm:faithful}).
\end{theorem}

\begin{proof}
\textup{(i)} is \Cref{def:ablation,thm:element-floor}. \textup{(ii)} a legal
reorder changes the interior path while fixing endpoints; invariance is
endpoint-determination---path opacity (the relaxation's landing depends on the
fixed point, not the route). \textup{(iii)} a same-slot synonym is a different
representation of the same element averaging to the same alignment; invariance
is representation-independence. \textup{(iv)} negation is a transport that, for
a faithful landing, must carry the target to its negation along a closing
cycle: correct inversion is $\hol=0$, incorrect $\hol\neq0$.
\end{proof}

\begin{remark}[Perturbation makes the floor a number]
\label{rem:perturb-floor}
One cannot read $\flo$ off a language; one can ablate a word and watch the
landing collapse, and the size of the collapse \emph{is} the item's
distance-to-floor. The suite turns the four structural facts (elements,
opacity, mobility, holonomy) into four empirical dials, and the four are the
faces of the one floor. This is the ``resolution validation by linguistic
perturbation'' of an earlier specification, now seen to measure the floor, not
four ad-hoc properties. No theorem depends on the reading.
\end{remark}

% =====================================================================
\section{The complete pipeline}
\label{sec:pipeline}
% =====================================================================

\begin{construction}[Translation pipeline]
\label{con:pipeline}
Given source $\Lang$, target $\Lang'$, and a source text:
\begin{enumerate}[label=\textup{(\arabic*)},leftmargin=2.6em]
\item \textbf{Terminate.} Cast each unit in completed aspect
(\Cref{con:perfective}); only terminated units propagate
(\Cref{thm:terminate}).
\item \textbf{Seed two columns.} Source as $\colA$; naive cognate translation
as $\colB$, the seed (\Cref{con:columns}).
\item \textbf{Anchor.} Ablation-rank each column; fix high-$\abl$ elements as
anchors, leave the residue free (\Cref{con:columns},
\Cref{thm:element-floor}).
\item \textbf{Relax.} Alternate $\relax_{\colA\to\colB}$ and
$\relax_{\colB\to\colA}$ (\Cref{def:relaxation}): perturb a column,
reverse-translate, propagate the demanded residue change back, iterate.
\item \textbf{Halt at quiescence---or decline.} If cross-demands reach zero,
the pair is quiescent: a proper translation, identical in meaning
(\Cref{thm:quiescence-identity}), at residual $\res\ge\flo$ of form
(\Cref{thm:form-residual}). If they do not vanish, the unit is deeply
untranslatable and the method declines (\Cref{cor:decline}).
\item \textbf{Resolve shallow untranslatables.} For a quiescent pair with no
single floor-distance target element, render by minimal composite or borrow
(\Cref{thm:two-kinds}(i)).
\item \textbf{Audit and validate.} Confirm faithfulness by back-cycle closure
(\Cref{thm:faithful}); flag false friends (\Cref{thm:falsefriend}); run the
perturbation suite (\Cref{def:perturb}) whose signature
(\Cref{thm:perturb}) calibrates confidence.
\item \textbf{Represent.} Choose a representation (register, style) of each
landed element (\Cref{thm:perturb}(iii)) and restore natural target aspect.
\end{enumerate}
The output, when it exists, is a target text in semantic quiescence with the
source---identical in meaning, certified blind by the vanishing of
cross-demands---with its load-bearing elements anchored and only its
form-residue distinguishing it.
\end{construction}

\begin{remark}[The runtime an earlier specification anticipated]
\label{rem:runtime}
This pipeline is the runtime a prior points-and-resolutions specification
reached for and could not ground. Its \emph{points} (content with inherent
uncertainty, never fully certain) are region-valued cells; its
\emph{resolutions} are residue propagations; its
\emph{remaining\_uncertainty} is the residual $\res\ge\flo$; its
\emph{perturbation validation} is \Cref{sec:faithful}; its \emph{fuzzy units}
of emergent grain are the elements of \Cref{thm:catalogue}; its keyword
\texttt{propagate\_uncertainty} is the propagation itself. It reached for
probability, fuzzy membership, and Bayesian updating because no floor had been
derived; here the single derived floor replaces all three, parameter-free.
The shape was right; it lacked its foundation. No theorem depends on the
reading.
\end{remark}

% =====================================================================
\section{The picture in one statement}
\label{sec:picture}
% =====================================================================

A text exists because uncertainty propagated; its meaning is what that
propagation emits, not its content, and content---which precision targets---is
the non-meaning. Meaning cannot be extracted, so the translator does not try
to know it; it places the source and a naive cognate seed as two columns,
anchors the unreplaceable elements of each by ablation, and relaxes: perturb
a column, reverse-translate, propagate the demanded change back, iterate.
When the relaxation halts---when a change on either side propagates and
demands nothing of the other---the two sides have the same meaning,
necessarily, because any difference of meaning is itself a demand and
quiescence is the absence of demands; and this identity is certified without
the meaning ever being read, because identity is no-difference and
no-difference is testable blind. The halt is at positive residual: same
meaning, irreducibly different form. When the relaxation does not halt, the
meanings cannot be made identical---deep untranslatability---and the honest
output is to decline. Given a grammatically wrong statement, the human
translates it (depending on meaning), the substituter mangles it (depending on
form), and our method translates it iff a fixed point still exists and
declines otherwise (depending on neither). This is translation without
meaning: the meaning-dependent translator's result, by a meaning-blind,
comprehension-free mechanism, sufficient exactly where a fixed point exists.

% =====================================================================
\section{Conclusion}
\label{sec:conclusion}
% =====================================================================

We have formalised translation, completely and on finite weighted graphs, as
the bidirectional relaxation of two columns to semantic quiescence. The
object of translation is the meaning that propagates from the source as a
cause, not its content---precision targets content and so targets the
non-meaning, as a precise spectrum is not a molecule's significance. Meaning
is an unextractable invariant, so the translator does not know it; it detects
that two sides \emph{share} it, by the absence of any propagable difference.
The elements held fixed in the relaxation are the ablation-unreplaceable
items, of emergent grain; the residue is what moves. The central theorem is
that quiescence is identity of meaning---necessarily, and certified blind,
because a meaning-difference is a demand and quiescence is no demand---reached
at positive residual of form, never at zero. Non-halting is the proof of deep
untranslatability, where the honest output is to decline rather than
hallucinate. A single dissociating experiment---a grammatically wrong
statement---separates the three translators by their dependency: the human on
meaning, the substituter on form, and our method on the existence of a fixed
point, reaching the human's result without comprehension and declining at the
boundary where the human infers and the substituter hallucinates. Faithfulness
is zero back-holonomy, implied by quiescence; the false friend is the bridge
caught only by route-audit; the perturbation suite measures the four faces of
the one floor. Throughout, the derived, parameter-free floor does the work of
every borrowed uncertainty calculus. This is translation without meaning:
sufficient for a translator that does not need to understand, because identity
of meaning is the silence of all demands between two sides, and that silence
can be reached, and recognised, blind.

\bibliographystyle{plainnat}
\begin{thebibliography}{9}
\bibitem[Banach(1922)]{banach1922} Stefan Banach. Sur les op\'erations dans
les ensembles abstraits et leur application aux \'equations int\'egrales.
\emph{Fundamenta Mathematicae}, 3:133--181, 1922.
\bibitem[Diestel(2017)]{diestel2017} Reinhard Diestel.
\emph{Graph Theory}, 5th edition. Springer GTM 173, 2017.
\bibitem[Ford and Fulkerson(1956)]{fordfulkerson1956} L.~R. Ford and
D.~R. Fulkerson. Maximal flow through a network.
\emph{Canadian Journal of Mathematics}, 8:399--404, 1956.
\bibitem[Menger(1927)]{menger1927} Karl Menger. Zur allgemeinen
Kurventheorie. \emph{Fundamenta Mathematicae}, 10:96--115, 1927.
\bibitem[Shannon(1948)]{shannon1948} C.~E. Shannon. A mathematical theory of
communication. \emph{Bell System Technical Journal}, 27:379--423, 1948.
\bibitem[Jakobson(1959)]{jakobson1959} Roman Jakobson. On linguistic aspects
of translation. In \emph{On Translation}, Harvard University Press, 1959.
\bibitem[Nida(1964)]{nida1964} Eugene A. Nida. \emph{Toward a Science of
Translating}. Brill, 1964.
\bibitem[Zadeh(1965)]{zadeh1965} L.~A. Zadeh. Fuzzy sets.
\emph{Information and Control}, 8(3):338--353, 1965.
\end{thebibliography}

\end{document}
